\section{Discussion}

\label{sec:results-overall-interpretation}

Taken together, the results show that weakening instrument relevance changes each dimension of finite-sample performance, but not in the same way. Point accuracy deteriorates systematically, accepted sets expand, and power against false candidate values generally declines, subject to the documented small reversals. Coverage is less systematic and can move toward its nominal level precisely when the procedure rejects few candidates. Near-nominal coverage can therefore coexist with low power and acceptance of most or all of the numerical domain. Accuracy, calibration, power, and informativeness measure different features of the procedure, and no one of them is sufficient to assess inference.

The estimator comparison reinforces this interpretation. Oracle-GMM knows the designated relevant controls, yet its point accuracy, power, and accepted-set localisation still deteriorate as relevance declines. High-dimensional nuisance estimation therefore cannot be the sole explanation for the common relevance pattern. DML-IVQR usually has lower RMSE than Full-GMM, but the documented exceptions and the separate coverage and power results rule out uniform superiority. Oracle is likewise an infeasible benchmark rather than a guaranteed finite-sample winner. Because the implementations differ in nuisance dimension, regularisation, and residualisation details, these are comparisons among complete estimator implementations; they do not isolate the causal effect of regularisation.

Sample size and quantile comparisons require similar qualifications. Increasing $n$ from 500 to 1000 lowers RMSE in all 60 matched comparisons and typically narrows accepted sets and raises power where usable signal remains. It does not improve every inferential dimension: several strong-relevance DML coverage rates deteriorate substantially, and power has documented exceptions. Across quantiles, the lower-tail biases are uniformly negative and the upper-tail biases uniformly positive under the convention $\mathrm{Bias}=\alpha_0-\widehat\alpha$. At fixed estimator, $n$, and $\kappa$, both MAE and RMSE are lower at the median than at either tail throughout the final results, whereas coverage has no corresponding universal centre-versus-tail ordering. Results at two sample sizes establish finite-sample contrasts, not an asymptotic rate or limit.

Neyman orthogonality is designed to reduce first-order sensitivity to small nuisance-estimation errors \parencite{chernozhukov_et_al_2018}, while candidate-specific test inversion is designed to avoid relying on a precisely estimated structural coefficient under weak identification \parencite{chernozhukov_hansen_2008}. Neither property creates identifying variation. Subject to its assumptions and regularity conditions, weak-identification-robust construction concerns validity of testing a candidate; it does not promise a narrow confidence region or high power when many false candidates are observationally difficult to distinguish. The DML-IVQR framework of \textcite{chen_huang_tien_2021} combines high-dimensional nuisance estimation with this style of candidate-specific inference, but the present simulations evaluate the preserved executable procedure. They do not validate the ideal asymptotic theory.

The population density weight at the true parameter encounters the non-finite regularity problem derived in Section~\ref{sec:framework}. The executable estimators instead use the fitted-normal density proxy and implementation-specific residualisation rules in Section~\ref{sec:implementation}; Appendix~\ref{app:reproducibility} further records that the preserved \texttt{dnorm} call supplies the residual variance in the argument position interpreted by R as a standard deviation. DML uses the same sample for nuisances and scores rather than cross-fitting, and its pivotal penalty generates fresh draws across evaluations. All three estimators use an asymptotic chi-square reference cutoff, while their empirical sizes differ. These features qualify the relationship between the experiment and the theoretical procedures, but the design does not vary them separately. None can therefore be identified as the cause of a particular distortion, and the results do not establish that changing one feature would remove it.

The numerical and design scope imposes further limits. Point estimates minimize a 41-point profile and hence reflect its resolution, endpoints, and first-tie rule. The accepted-set measure is restricted to $A_0=[-1,3]$ and cannot describe acceptance outside that domain; Appendix~\ref{app:numerical-domain-diagnostics} shows that accepted measure and boundary acceptance can continue to grow on wider investigated domains under weak relevance. Those diagnostics establish limited numerical localisation over the ranges examined, not global unboundedness or non-identification. More broadly, the evidence comes from one DGP family, five quantiles, two sample sizes, four relevance settings, and three preserved estimator implementations. It cannot establish how every alternative DGP, implementation, or larger sample would behave.

The bounded implication is therefore to distinguish what can be assessed in a simulation from what can be established in an empirical application. A procedure can produce a more accurate point estimate while remaining poorly calibrated, attain coverage near the nominal level while excluding few false candidates, or show a finite-domain measure near its ceiling without establishing the shape of an unrestricted confidence region. Simulation evaluations should report accuracy, calibration, power, and informativeness jointly. In empirical applications, researchers should disclose numerical-domain sensitivity and the assumptions underlying inference; nominal confidence levels alone do not establish finite-sample reliability. Within the scope of this experiment, that distinction gives the clearest account of what is lost as instrument relevance weakens.
